\section{System Architecture}
\label{sec:architecture}

\subsection{Pipeline Overview}

Figure~\ref{fig:pipeline} shows the pipeline's five stages. A sample
is analyzed statically and dynamically in parallel; both outputs feed
a \emph{normalizer} that deduplicates IOCs and groups ATT\&CK
observations by technique; a \emph{mapper} reconciles cross-source
confidence and produces the final technique list; an \emph{exporter}
converts that list into a STIX~2.1 bundle and/or a live MISP event.
Every stage is a standalone command-line tool operating on a
well-defined JSON schema, so the pipeline can be run end to end or any
single stage re-run in isolation, a property used extensively during
validation: a single detection-rule change often only requires
re-running the mapper, not the full sandbox detonation.

\begin{center}
\begin{tabular}{c}
\texttt{sample} \\
$\downarrow$ \\
\texttt{Static Analysis} \quad + \quad \texttt{Dynamic Analysis (CAPE)} \\
$\downarrow$ \\
\texttt{Normalizer} $\rightarrow$ \texttt{normalized\_iocs.json} \\
$\downarrow$ \\
\texttt{Mapper (confidence reconciliation)} $\rightarrow$ \texttt{attck\_mapping.json} \\
$\downarrow$ \\
\texttt{Exporter} $\rightarrow$ \texttt{STIX 2.1 bundle} / \texttt{MISP event}
\end{tabular}
\captionof{figure}{Pipeline stages and intermediate artifacts.}
\label{fig:pipeline}
\end{center}

\subsection{Static Analysis Engine}
\label{sec:static-rules}

The static engine parses PE (native and .NET) and ELF binaries,
extracting imports, strings (standard and FLOSS-decoded stack
strings~\cite{floss}), sections, YARA matches, and entropy
characteristics. For installer-packaged samples (NSIS, Inno, MSI) it
performs multi-pass extraction: every bundled file is extracted, keys
are discovered across \emph{all} extracted files before any decryption
is attempted (so a key embedded in one component can decrypt a payload
in another), candidate keys are tried against any file that looks
encrypted, and every file, decrypted or not, is then fully analyzed.
\S\ref{sec:case-installer} details a real bug found and fixed in how
this aggregation used to work.

Detection rules map static evidence to ATT\&CK techniques through
three complementary mechanisms:

\begin{itemize}
  \item \textbf{Import/string tables.} A curated mapping from specific
    Windows API names or string patterns to technique identifiers.
  \item \textbf{Combination checks.} For techniques whose single
    strongest indicator is also common in benign software, a rule
    requires a specific \emph{coherent combination} of indicators
    before reporting elevated confidence. Table~\ref{tab:combos} lists
    every combination check implemented via explicit multi-import
    matching; the same corroboration-count mechanism also gates several
    string-pattern rules added during coverage extension (e.g.\ Mshta's
    \code{mshta.exe} + \code{.hta}, WMI Event Subscription's three
    corroborating class names), tabulated in Appendix~\ref{app:rules}
    instead.
  \item \textbf{.NET BCL calls.} A .NET binary's native import table is
    typically only the CLR bootstrap stub, so managed method bodies are
    scanned directly for calls to specific Base Class Library APIs
    (e.g.\ \code{Convert::FromBase64String}).
\end{itemize}

\begin{table}[h]
\centering
\scriptsize
\begin{tabular}{@{}p{3.0cm}p{6.6cm}p{3.0cm}@{}}
\toprule
\textbf{Technique} & \textbf{Required combination} & \textbf{Rejected alone} \\
\midrule
T1055 (Process Injection) & \code{VirtualAllocEx} + \code{WriteProcessMemory} + \code{CreateRemoteThread}, or the read-primitive trio \code{OpenProcess} + \code{VirtualQueryEx} + \code{ReadProcessMemory} & \code{OpenProcess} alone \\
\rowline
T1055.012 (Process Hollowing) & \code{CreateProcessW} + \code{WriteProcessMemory} + \code{SetThreadContext} + \code{ResumeThread} + \code{Nt/ZwUnmapViewOfSection} & any single API \\
\rowline
T1685 (Disable or Modify Tools) & \code{WTSEnumerateProcessesW} + \code{OpenProcess} + \code{TerminateProcess} & bare \code{TerminateProcess} \\
\rowline
T1010 (Application Window Discovery) & \code{EnumWindows} + \code{GetWindowTextW} & bare \code{GetWindowTextW} \\
\rowline
T1134.001/.003 (Impersonation/Theft; Make and Impersonate Token) & \code{ImpersonateLoggedOnUser} + \code{DuplicateTokenEx} (.001) or \code{LogonUserW} (.003) & bare \code{Impersonate\-LoggedOnUser} \\
\bottomrule
\end{tabular}
\caption{Combination-aware detection rules. A component listed as
``rejected alone'' is intentionally excluded from single-import
confidence promotion because it has substantial legitimate use outside
the claimed technique.}
\label{tab:combos}
\end{table}

\subsection{Dynamic Analysis Engine}
\label{sec:dynamic}

Dynamic detonation runs inside CAPE~\cite{cape} against an isolated,
INetSim-simulated network with no real Internet egress, a
containment-first design whose consequence for observability is
discussed in \S\ref{sec:limitations}. CAPE's raw report is curated
into a smaller, IOC-focused JSON: signatures, dropped files, network
activity, and a curated ATT\&CK mapping list.

CAPE's community signatures carry their own technique tags, but these
were treated as a hypothesis, not ground truth, throughout validation.
The curation layer maintains three correction tables, each entry
individually justified against the signature's own raw evidence (its
\code{data} field, not just its name or category):

\begin{itemize}
  \item \textbf{Unreliable signatures}, dropped entirely: e.g.\
    \code{unbacked\_process\_mitigation\_alteration} (tags
    \technique{T1562}), whose raw caller addresses across the sample
    that triggered it all sit in a narrow high-memory range consistent
    with a system DLL, suggesting the ``unbacked memory'' classification
    misfired on Windows' own automatic process initialization rather
    than the sample's own action.
  \item \textbf{Per-signature technique drops}, where one tag on an
    otherwise-reliable signature is wrong: e.g.\
    \code{antiav\_servicestop} tags \technique{T1489} (Service Stop),
    but on the sample that triggered it, its own raw data names the
    specific service stopped: the malware's \emph{own} driver service,
    not a security product's.
  \item \textbf{Technique remaps}, where MITRE's own definitions
    contradict the community tag: e.g.\ \code{unbacked\_api\_resolution}
    is community-tagged \technique{T1129} (Shared Modules), which MITRE's
    own page defines as requiring a module backed by a file on disk;
    the signature's own description (``dynamically allocated,
    unbacked memory'') is the textbook opposite, mapped instead to
    \technique{T1620} (Reflective Code Loading). The same table also
    bridges a version gap: CAPE's community signatures are frozen to
    whichever ATT\&CK version they were last written against, so
    \technique{T1562}/\technique{T1562.001} and \technique{T1070.001}
    tags fired by CAPE are remapped to their ATT\&CK v19 identities
    (\technique{T1685}, \technique{T1685.005}) at curation time,
    independent of this project's own pinned version
    (\S\ref{sec:limitations}).
\end{itemize}

A further capability, added during coverage extension, checks CAPE's
raw \code{executed\_commands} against known LOLBin/reconnaissance
command patterns that have no corresponding community signature at
all: e.g.\ \code{netsh wlan show profiles}, which CAPE observes and
records but does not itself map to any technique.

\subsection{Confidence Reconciliation}
\label{sec:reconciliation}

The core reconciliation rule (\S\ref{sec:reconciliation}, implemented
in \code{reconcile.py}) is: a technique observed by both static and
dynamic evidence reaches \emph{high} confidence unless both sources'
own best confidence is \emph{low}, in which case it is \emph{medium};
a technique observed by only one source keeps that source's own best
tier, and cannot be promoted by volume of evidence from a single
source alone.

\subsubsection{Cross-Level Extension}

This rule, as originally implemented, compared observations grouped by
\emph{exact} technique identifier. In practice, a dynamic signature
frequently reports a parent technique (e.g.\ \technique{T1547}) while a
static rule reports the specific sub-technique responsible
(\technique{T1547.001}): the same underlying persistence mechanism,
observed at two different levels of specificity by two different
sources, but never corroborating each other under exact-match
grouping. This is a real inconsistency, not a hypothetical one: the
evaluation harness's own family-level matching (\S\ref{sec:evaluation})
already treats a parent technique and its sub-technique as equivalent
for scoring purposes, while the confidence model that produces the
scored output did not apply the same principle.

\code{apply\_cross\_level\_corroboration()} closes this gap: after the
normal per-technique reconciliation pass, techniques are grouped by
base identifier (stripping any sub-technique suffix), and for any group
where static \emph{and} dynamic evidence exist somewhere among its
members but an individual member is missing one side, that member's
confidence is recomputed using the group's best tier for the missing
side, the same ``high unless both low'' rule, extended across sibling
identifiers rather than only within one. Every promotion is recorded
on the technique's own record (which sibling technique supplied the
missing corroboration), so the mechanism is auditable rather than
silently blending evidence across identifiers. On the three validated
samples this change did not alter any final confidence value: every
currently-mixed-granularity group either already independently reached
\emph{high} through its own source's severity, or both members shared
the same single source. It was nonetheless built \emph{before}, not
after, the coverage extension described next, specifically because new
sub-technique-specific static rules paired with CAPE's often
parent-level-only community signatures make this gap more likely to
matter going forward, not less.

\subsection{The Resubmission Loop}
\label{sec:resubmission}

Multi-stage malware's real capability is frequently distributed across
several binaries: RoningLoader alone, across its infection chain,
drops or reflectively loads a kernel rootkit driver, an antivirus-killer
DLL, and a separate remote-access-trojan C2 client
(\S\ref{sec:case-roning}). Naively merging every dropped component's
evidence into the parent sample's own \code{attck\_mapping.json} would
misattribute confidence: \code{reconcile.py}'s model assumes every
observation grouped under one technique describes the \emph{same}
binary, and a dropped payload's own capability is not proof the parent
loader itself performs that behavior.

The resubmission loop instead gives each dropped/extracted component
its own, fully independent run through the same static
analysis $\rightarrow$ normalize $\rightarrow$ map pipeline, tagged
with explicit lineage (parent hash, discovery mechanism, discovery
timestamp) back to the sample that produced it. Three components,
split by dependency rather than convenience:

\begin{enumerate}
  \item \textbf{Manifest writer} (host-run, pure standard library):
    locates dropped-file bytes in CAPE's own storage, copies them into
    a queue, and writes a lineage manifest per candidate, bounded by a
    configurable artifact cap.
  \item \textbf{Static processor} (runs inside the static analysis
    container, which already has the heavier dependencies, \code{pefile},
    \code{yara-python}, FLOSS, Ghidra, that dropped-file analysis
    needs): analyzes each queued artifact and tags its report with the
    lineage manifest's data. Duplicating those dependencies into a
    separate, lighter container was considered and rejected as
    unnecessary maintenance surface for no isolation benefit.
  \item \textbf{Merge processor} (runs inside the pipeline container,
    pure standard library): finds every lineage-tagged report not yet
    normalized and mapped, and runs it through the same
    confidence-reconciliation pipeline as a top-level sample, writing
    to a component-specific output path, never the parent sample's own
    \code{attck\_mapping.json}, so a dropped component's evidence can
    never silently overwrite or blend into the parent's own scored
    output. This enforces the same non-blending principle at the
    file-system level that \code{reconcile.py} enforces at the
    data-model level.
\end{enumerate}

For RoningLoader specifically, the resubmission loop discovered and
independently analyzed 26 components beyond the parent loader binary:
the rootkit driver, the AV-killer DLL, and the RAT client, each with
its own dedicated entry in the manual ground truth, plus 23 further
extracted payload and configuration artifacts with no individually
authored ground truth of their own. Section~\ref{sec:evaluation} scores
the three named components against their per-component ground truth,
since parent-only scoring structurally understates what the full
pipeline, resubmission loop included, found; the remaining 23 are
additional real evaluation surface in a different sense
(\S\ref{sec:limitations}, item~1), each confirmed to produce sane,
non-crashing, lineage-tagged output rather than contributing new scored
precision/recall data points.

\subsection{Detection Rule Coverage Extension}
\label{sec:coverage}

Every static rule present at the start of the validation described in
\S\ref{sec:evaluation} existed because specific evidence in one of the
three validated samples justified it, a sound approach for avoiding
overfitting, but one that leaves coverage narrow relative to the full
MITRE ATT\&CK matrix, not just relative to what the three validated
families individually need. To quantify this precisely rather than
estimate it, the current official MITRE Enterprise ATT\&CK STIX
bundle~\cite{mitre-attack-stix} was downloaded and counted directly:
474 Windows-relevant techniques (including sub-techniques) in the
present matrix, of which roughly 30 had a dedicated static rule before
this extension.

Coverage was subsequently extended across six batches, one to two
tactics at a time, adding approximately 55 further technique
identifiers (full list in Appendix~\ref{app:rules}). Every addition
follows the same sourcing discipline: the detection pattern is the
\emph{literal} mechanism MITRE's own technique description names, a
specific registry value, a WMI class combination, a file marker, a
DLL, or a command-line flag, never a generic API or tool name alone
when that name also has ordinary legitimate uses. After each batch,
static analysis was re-run on all three validated samples and checked
for any new finding before the batch was considered safe to keep;
none of the 55 additions fired on any of the three samples, confirming
the extension does not introduce new false-positive risk on the set
this project's own numbers are measured against. This is explicitly
coverage aimed at samples not yet seen; it does not, and is not
claimed to, move the evaluation numbers reported in
\S\ref{sec:evaluation}.

Coverage extension was scoped deliberately, not exhaustively, after
checking real technique counts for every remaining MITRE tactic
(\S\ref{sec:future-work}): Reconnaissance and Resource Development are
structurally out of scope regardless of effort, since both describe
attacker activity that occurs before the malware sample exists;
Initial Access and Lateral Movement are technically in scope but poor
fits for what single-binary detonation analysis can observe.

\subsection{Deployment and Reproducibility}
\label{sec:deployment}

The prototype and every analysis environment it depends on, not only
the pipeline's own code, is defined in a single \code{docker-compose.yml},
split into two profiles activated independently. The \code{core}
profile (static analysis engine, pipeline, ATT\&CK Navigator) needs no
host-level setup beyond Docker itself and is what produced every
static-only result in this thesis. The \code{sandbox} profile (CAPE,
MISP, and their supporting services: MySQL, Redis, INetSim)
additionally needs one-time host configuration of libvirt/KVM and a
snapshot-able Windows guest VM, since a VM disk image cannot itself be
provisioned by \code{docker-compose}. Table~\ref{tab:deployment} lists
every service and its profile.

\begin{table}[h]
\centering
\scriptsize
\begin{tabular}{@{}llp{6.6cm}c@{}}
\toprule
\textbf{Service} & \textbf{Profile} & \textbf{Role} & \textbf{Memory} \\
\midrule
\code{static} & core & PE/ELF static engine (\code{pefile}, YARA, FLOSS, Ghidra headless) & 3g \\
\rowline
\code{pipeline} & core & Normalizer, mapper, exporter (STIX/MISP) & 1g \\
\rowline
\code{navigator} & core & ATT\&CK Navigator layer visualization & 512m \\
\rowline
\code{redis} & sandbox & MISP session/job backend & 256m \\
\rowline
\code{inetsim} & sandbox & Simulated network for detonation, no real egress & 256m \\
\rowline
\code{cape} & sandbox & CAPE sandbox controller (KVM/libvirt detonation) & 6g \\
\rowline
\code{mysql} & sandbox & MISP's relational store & 1g \\
\rowline
\code{misp} & sandbox & MISP core (event storage, API, galaxy data) & 2g \\
\bottomrule
\end{tabular}
\caption{Service inventory by profile, from \code{docker-compose.yml}.}
\label{tab:deployment}
\end{table}

The one setup step \code{docker-compose} genuinely cannot perform is
provisioning the KVM/libvirt hypervisor and its virtual network on the
host itself, since \code{cape} and \code{inetsim} need a real
\code{/dev/kvm} and a real libvirt bridge to attach to, not a container
namespace. \code{scripts/host-prereqs.sh} automates this one-time step
and is idempotent: it installs \code{qemu-kvm} and \code{libvirt},
verifies \code{/dev/kvm} is actually usable, and discovers the bridge
interface, subnet, and gateway IP libvirt is using on that specific
machine (deliberately not hardcoded, since it varies host to host),
writing the discovered values into \code{.env} so \code{inetsim} and
CAPE's own config interpolation can pick them up without an image
rebuild.

Two further design decisions in the compose file follow directly from the
containment-first posture described in \S\ref{sec:dynamic}. First,
backend services (static analysis, pipeline, MySQL, Redis) sit on a
Docker network with \code{internal: true} and no route to the host or
the internet at all; only MISP and the ATT\&CK Navigator, the two
services a human needs to reach from a browser, sit on a second,
non-internal network, keeping the no-egress boundary at the narrowest
point that still lets a human review results. Second, CAPE and INetSim
run with host networking rather than on the isolated backend network, a
deliberate, scoped exception: CAPE's own KVM/libvirt-based detonation
has no officially supported Docker deployment pattern, and every
community approach converges on host networking for the sandbox
controller itself, so the no-egress boundary is enforced one layer
down, at the guest VM's own network configuration, rather than at the
Docker layer for this one service. Running both profiles together on a
16GB host budgets to roughly 14GB, which only fits because detonation
is sequential, one guest VM at a time, never concurrent.
